\section{Source Index Contract}

Source hashes bind raw bytes. Normalized hashes bind canonical records. The
source index is the bridge between those layers.

\begin{center}
\small
\begin{tabularx}{\textwidth}{lX}
\toprule
Field & Contract \\
\midrule
\texttt{source\_id} & Stable package identifier for the source artifact. \\
\texttt{path} & Package-relative path under \texttt{sources/}; absolute paths and parent-directory escapes are invalid. \\
\texttt{sha256} & Domain-separated hash over the raw source bytes. \\
\bottomrule
\end{tabularx}
\end{center}

The current synthetic source index is intentionally small:
\begin{quote}
\footnotesize
\begin{verbatim}
{
  "source_id": "source:synthetic-summary-export",
  "path": "sources/synthetic-summary-export.json",
  "sha256": "sha256:f1fd7040..."
}
\end{verbatim}
\end{quote}

Real adapters will need more metadata: source format, export system, parser
version, public-record citation, custody note, redaction status, and parser
diagnostics. That richer adapter contract belongs to V.9. V.3 only requires the
baseline replay rule: the package names the source path, hashes the raw bytes,
and rejects empty or missing source evidence.
